\chapter{Background and related work}
\label{ch:background}

Several literatures bear on the claim this dissertation tests, and none of them
measures it. Scaling laws describe loss as a function of parameters, data and
compute, but not what a fixed checkpoint can decide. Tokenisation-free modelling
explains why a byte-level encoder is possible, not whether its decisions are
correct. The long-context literature measures how much of a context a model uses,
not whether a constrained selector survives a wider candidate set. Work on
ranking, constrained decoding, set-valued prediction and selective classification
supplies the formal vocabulary for selection among candidates, but it optimises
different objectives: ordering quality, surface-form validity, coverage, or
risk--coverage, rather than the accuracy of a single choice as the bank widens.
This chapter separates what each literature establishes from what it leaves open,
and states where \textsc{Azurite} sits.

\section{Scaling laws and their limits}

\citet{kaplan2020scaling} established that language-model cross-entropy falls as
a power law in parameters, data and compute when all three move together. The
result is a family of curves; it contains no statement about what happens when
parameters move while data is held fixed, which is the slice most small-scale
architecture comparisons actually take. \citet{hoffmann2022chinchilla} closed the
compute-optimal case: at a fixed budget, parameters and tokens should scale
together at roughly twenty tokens per parameter. A model at $\DN = 0.88$ is about
$23\times$ below the token count its size calls for, so its poor held-out score is
the predicted outcome rather than an anomaly.

The Chinchilla result has itself been replicated and qualified.
\citet{besiroglu2024replication} reproduce the fit and report that the estimated
exponents and the implied optimal allocation depend on the fitting procedure and
the data retained; the paper is a warning about treating a three-point fit as a
law. \citet{muennighoff2023data} study the data-constrained regime and find that
repeating epochs is close to fresh data for a limited number of repetitions,
after which returns decay sharply. \citet{sardana2023beyond} show that once
inference cost is included, the loss-optimal allocation shifts toward smaller
models trained on more tokens. All three matter here: the ladder is deliberately
below the compute-optimal line, so its loss is not a claim about what the
architecture would achieve at $20$ tokens per parameter, and the suite's fixed
token budget places different rungs at different distances from that line.

The inference from loss to decisions is the weakest link. \citet{schaeffer2023emergent}
show that some ``emergent'' ability curves are artifacts of the metric: a
discontinuous or thresholded measure can manufacture a jump that a continuous
loss does not contain. The lesson generalises beyond emergence. Held-out bits per
byte and task accuracy are different measurements, and a clean loss law does not
imply a clean accuracy law. This dissertation reports both and treats their
divergence as a result rather than an inconsistency.

\section{Memorisation, double descent and the objective}

\citet{zhang2017rethinking} showed that networks fit random labels perfectly while
still generalising on real ones: the capacity to fit is not the capacity to
generalise. In the data-limited direction the same capacity explains the large
train/held-out gap that the pre-loss sweep reported: when a model has more
parameters than the corpus has bytes, storing the corpus is the cheapest available
descent direction. That sweep is not available in this checkout, so the claim is
stated as a mechanism, not as a surviving measurement.

The resulting loss curve---falling, then rising, as capacity grows---resembles
double descent \citep{belkin2019reconciling,nakkiran2020deep}. This dissertation
does not claim to have demonstrated double descent. The surviving sweep is too
incomplete to test the shape: the mid and large cells were reclaimed. The
memorisation account makes a prediction that can be tested (the epoch sweep in
\S\ref{sec:e2}); ``double descent'' here would name the shape rather than explain
it.

The objective interacts with that boundary. \citet{devlin2019bert} used $15\%$
masking, tuned for a small token budget; \citet{portes2023mosaicbert} report that
$30\%$ masking trains roughly twice as fast at equal compute, because more
positions carry gradient per forward pass. \textsc{Azurite} uses $30\%$. At that
rate and $\DN < 1$ the model sees each byte roughly once per epoch as a target,
which is part of why the memorisation boundary sits near $\DN \approx 1$--$2$
rather than at $\DN \approx 0$.

\section{Tokenisation and byte-level modelling}

Subword tokenisation is the default because it shortens sequences.
\citet{sennrich2016bpe} introduced byte-pair encoding to neural machine translation to handle rare words,
and \citet{kudo2018sentencepiece} made a language-independent implementation
practical. The cost is a fixed vocabulary and a preprocessing boundary: the model
cannot represent a string the tokeniser has not assigned an id, and the mapping
from bytes to tokens is outside the learned parameters.

Token-free models remove that boundary. \citet{xue2022byt5} show that a
byte-to-byte encoder is competitive with subword models and more robust to noise,
at the cost of longer sequences. \citet{clark2022canine} and
\citet{tay2021charformer} avoid subword vocabularies but reintroduce a learned
downsampling stage, which is a tokeniser in all but name. \citet{yu2023megabyte}
takes the more ambitious route: fixed-size byte patches processed by a
multi-scale transformer, so that sequence length grows with patching rather than
with bytes. \citet{pagnoni2024blt} replace fixed patches with dynamically chosen
ones and argue that patches, not tokens, are the right scaling unit.

\textsc{Azurite} takes the byte-level route for two reasons. There is no tokeniser
boundary, so text, code and quantised audio can share one input space; and in a
decider the expensive operation is encoding the context once rather than decoding
many bytes. The cost is real: at four times the sequence length of a four-byte
tokeniser, attention is sixteen times more expensive per unit of text. The claim
therefore has to be stated as accuracy per FLOP, not accuracy alone, and
Chapter~\ref{ch:results} reports both.

What the byte-level literature does not provide is a decision result. ByT5 and
CANINE are evaluated on generation and classification with free outputs; none of
them trains a selector at several candidate widths on one checkpoint. The
architectural argument for bytes is a statement about what the interface permits,
not evidence that a byte-level selector is accurate at width.

\section{Long context: mechanisms and measurement}

The transformer's attention is quadratic in sequence length \citep{vaswani2017attention},
which is why long-context work is mostly about changing the cost or the position
encoding. Rotary position embeddings \citep{su2021roformer} make relative position
a rotation and are the base on which later extension methods operate.
\citet{chen2023pi} extend a trained window by interpolating position indices, and
\citet{peng2023yarn} add a frequency-dependent scaling with a temperature-like
correction. \citet{dao2022flashattention} reduces the memory traffic of exact
attention without changing its mathematics; \citet{beltagy2020longformer} sparsify
attention for long documents. \citet{xiao2023streaming} show that a small set of
``attention sink'' tokens carries much of the stability of streaming inference,
and \citet{gu2023mamba} and \citet{poli2023hyena} offer subquadratic sequence
models as alternatives to attention altogether.

Mechanism is not capability. \citet{liu2023lost} find that retrieval accuracy is
highest when the relevant passage is at the beginning or end of the context and
degrades in the middle. \citet{hsieh2024ruler} and \citet{bai2023longbench}
construct tasks whose difficulty grows with length and report that models'
effective context is often far shorter than the advertised window. A model that
accepts $32$k tokens has not been shown to use them.

\textsc{Azurite}'s measured context is one point: a $256$-token sequence scores
$0.9901$ bits/byte, and the $1{,}024$- and $4{,}096$-token runs were reclaimed
with the first sandbox (Chapter~\ref{ch:results}). The long-context claim in the
product framing is therefore a design target, not a measured property. The
literature above explains why accepting a longer window is not the same as using
it, and why a decider that encodes a long context once still has to demonstrate
that the relevant evidence survives the encoding.

\section{In-context learning and task induction}

\citet{brown2020gpt3} established that a language model can perform a task from a
few demonstrations without weight updates. \citet{wei2022cot} showed that
intermediate reasoning steps improve multi-step tasks at scale.
\citet{olsson2022induction} located a mechanism---induction heads---that copies
and continues patterns seen in context, and \citet{wei2023larger} report that
larger models can override semantic priors from demonstrations in favour of
abstract task structure. \citet{akyurek2022what} and \citet{xie2022explanation}
give complementary accounts of in-context learning as an implicit algorithm or
Bayesian inference, and \citet{min2022rethinking} show that the labels in
demonstrations matter less than the format and the label space.

The relevance to a decider is indirect but real. A decider does not prompt and
does not update weights at inference; it fine-tunes a selector on a small set of
labelled rows and then scores a finite bank. But the in-context results explain
why the \emph{label space} is part of the task: changing the candidate set changes
what the model can infer about the task, even when the context is unchanged. That
is the mechanism the width sweep manipulates, and it is why the sweep trains the
width rather than restricting a fixed model. The in-context literature does not
measure that manipulation; it studies prompts, not candidate banks.

\section{Selecting among candidates}

\subsection{Ranking and reranking}

Learning to rank formalises selection over a query-conditioned set. Pointwise
methods score each item independently, pairwise methods learn relative preferences
\citep{burges2005ranknet}, and listwise methods optimise a distribution over the
whole list \citep{cao2007listnet}; \citet{liu2009ltr} reviews the three families.
Reranking applies the idea to an existing n-best list: \citet{shen2004reranking}
score translation hypotheses with a discriminative model, and
\citet{lee2021reranking} do the same for neural machine translation.

\textsc{Azurite}'s scorer is pointwise: each candidate is embedded independently
and scored against the context, with no cross-candidate attention. That is a
deliberate efficiency choice---one context encoding, $K$ candidate encodings, no
$K^2$ interaction---and it has a known cost. A listwise reranker can exploit
contrast between candidates; a pointwise scorer cannot. The width sweep measures
the consequence of width for this pointwise design, not for a listwise one.

\subsection{Constrained decoding is not selection}

A separate line restricts generation to strings that satisfy a constraint.
\citet{hokamp2017grid} use grid beam search to require phrases,
\citet{post2018fast} allocate beams dynamically to lexical constraints,
\citet{scholak2021picard} parse incrementally to constrain SQL output, and
\citet{lu2020neurologic} add predicate-logic constraints with lookahead;
\citet{geng2023grammar} constrain by a grammar and report strong structured-output
results without fine-tuning. Contrastive decoding \citep{li2022contrastive} and
contrastive search \citep{su2022contrastive} shape the distribution rather than
impose a hard constraint.

These methods guarantee surface form: the output satisfies the constraint, or it
is rejected. They do not guarantee that the constrained string is the correct
answer, and most of them do not operate on a declared candidate set at all. A
decider's structural guarantee is stronger in one direction---no out-of-bank
label---and weaker in another: it must rank the bank correctly, and the
constrained-decoding literature offers no measurement of that as the bank grows.

\subsection{Set-valued prediction, conformal sets and selective classification}

Set-valued classification returns a subset of labels rather than one. The
objective is a trade-off between coverage and set size:
\citet{sadinle2019least} construct least-ambiguous sets with bounded error,
\citet{mortier2021efficient} optimise a decision-theoretic utility over sets, and
\citet{chzhen2021set} give a unified framework. Conformal prediction supplies
finite-sample marginal coverage under exchangeability; \citet{romano2020classification}
introduce adaptive sets, \citet{angelopoulos2021uncertainty} regularise their size,
and \citet{angelopoulos2023gentle} is the standard introduction.
\citet{bates2021risk} extend the guarantee from coverage to a risk bound.

These are different claims from the one this thesis tests. A conformal set is
allowed to contain several labels and is judged on coverage and size; a decider
must return one label and is judged on accuracy. Coverage does not imply that the
top-ranked candidate is correct, and a set-valued method does not train a head at
a fixed width. The vocabulary is nevertheless useful: the width sweep is, in
effect, a study of how the accuracy of a single choice changes as the label set
grows, which set-valued work holds fixed by construction.

Selective classification is closer in spirit. \citet{geifman2017selective} trade
coverage for risk by abstaining on low-confidence examples, and
\citet{geifman2019selectivenet} learn the reject option end to end. Calibration is
the prerequisite: \citet{guo2017calibration} show that modern networks are poorly
calibrated and that temperature scaling helps, \citet{nixon2019measuring} catalogue
the pitfalls of expected calibration error, and \citet{kumar2019verified} give
estimators with explicit guarantees. The suite reports ECE and a risk--coverage
curve for exactly this reason; it does not claim an abstention guarantee, because
the curve is measured, not certified.

\subsection{Set architectures}

Permutation-invariant set encoders provide the architectural alternative to
pointwise scoring. \citet{zaheer2017deep} show that a sum-decomposed representation
is sufficient for permutation-invariant functions, \citet{lee2018set} add
attention between set elements, \citet{carion2020detr} use a set loss for object
detection, and \citet{vinyals2015pointer} point to elements of an input set.
\textsc{Azurite} uses none of these: candidates are scored independently. A
set-attention decider would condition each candidate on the others and could, in
principle, represent contrast; it would also pay $K^2$ attention per decision and
would make the width sweep a different experiment. That trade is stated here
rather than assumed away.

\subsection{Candidate generation and retrieval}

At scale, the selector is only as good as the bank it is given. Dense retrieval
\citep{karpukhin2020dpr}, retrieval-augmented generation \citep{lewis2020rag},
REALM \citep{guu2020realm} and FiD \citep{izacard2020fid} all separate candidate
generation from answer production. Their recall bounds the final accuracy: a
selector cannot choose a candidate that was never retrieved. The suite avoids
this confound by using a closed label bank per task, so every gold label is in the
bank by construction. That makes the width experiment interpretable, and it also
means the results say nothing about open-vocabulary retrieval recall.

\subsection{The decider claim and its public evidence}

The commercial decider is TypeSafe's Jev. Its launch post claims a parallel
sampling mechanism over a pre-declared output set, a price of \$0.042 per million
input tokens with output billed at zero, and a 70--500\,ms response band; it
publishes no parameter count, architecture, weights or per-task accuracy
\citep{typesafe2026jev}. The Laya model card describes a fine-tuned ``System 1
decision model'' but, as checked for this thesis, states no publication date and
no peer-reviewed result \citep{laya2026card}. The third-party evidence is uneven:
\citet{ayautomate2026jev} run Jev 1.13 through OpenRouter on 791 labelled
decisions, including a 77-way Banking77 cell where Jev scores 78.8\% (95\% Wilson
73.1--83.6); \citet{abdelstark2026jev} maintain a probability-aware evaluation
harness; and \citet{bonn2026jev} evaluate Jev zero-shot across 37 datasets in an
arXiv preprint. A peer-reviewed applied paper integrates Jev into a manufacturing
proof of concept \citep{vodilka2026jev}, but it is not a benchmark.

Two consequences follow. First, every task-level Jev number available to this
thesis comes from a harness it did not run: the third-party cells fix the rows and
the option list but not the hardware, and the vendor's own table does not document
its splits. A like-for-like comparison therefore requires running both systems on
the same rows, which this dissertation records as out of scope. Second, what can
be reported is a point-estimate comparison with its evidence class attached, not a
paired win. At 113.47M, \textsc{Azurite} clears the cited point on \SuiteLargeClear{} of
\SuiteComparatorCells{} comparator cells --- Banking77 \SuiteLargeBank\% against
78.8\%, Financial PhraseBank, DAIR Emotion and MASSIVE intent --- ties CLINC150,
and is below the cited point on AG News, XNLI and Language ID
(Chapter~\ref{ch:comparator}). The stored record therefore supports clear point
wins, including the pre-registered 77-way Jev cell, and a capability boundary on
the other cells; it does not support a blanket benchmark victory. The Banking77
win is against the pre-registered 78.8\% zero-shot comparator; the strongest
public Jev run reports 92.40\% on the full 3,080-row test, above
\SuiteLargeBank\%, so the flagship win does not survive that comparator
\citep{smith2026jevbanking}. The protocol-matched test on those rows and the
same 24 examples returns \JevMatchPctAOne\% (\JevMatchAOneLo--\JevMatchAOneHi)
on the primary arm and \JevMatchPctNoEx\% (\JevMatchNoExLo--\JevMatchNoExHi)
without examples, so the gap is not a row-set artifact. A separately
pre-registered query-anchored head then lifts the example arm to
\QueryPctQueryTwo\% and confirms mean-pool dilution as the mechanism, but it
remains \QueryDiffAbs{} points below Jev, so the gap is an information-budget
limit rather than a pooling limit (Chapter~\ref{ch:comparator}). Of the four
pre-registered clear wins, two are against Jev points (DAIR Emotion and
Financial PhraseBank); MASSIVE intent's comparator is Laya, and Banking77 does
not survive the stronger Jev run.

\section{Multimodality and a shared byte space}

A byte-level input makes a shared representation possible, not actual. Speech
models such as wav2vec~2.0 \citep{baevski2020wav2vec} and Whisper
\citep{radford2022whisper} learn from raw or spectrogram audio, and neural audio
codecs such as EnCodec \citep{defossez2022encodec} quantise audio into discrete
units that a sequence model can consume. ImageBind \citep{girdhar2023imagebind}
learns a joint embedding across modalities; Chameleon \citep{chameleon2024} and
Transfusion \citep{zhou2024transfusion} train early-fusion models over discrete
and continuous tokens; Gemini \citep{gemini2023} is a large multimodal family.
None of this is evidence about \textsc{Azurite}: no multimodal data has been
trained on, and no multimodal result exists. The architectural argument is a
statement about what the interface permits. It becomes a claim only when a
multimodal corpus and an evaluation exist, and this dissertation does not supply
them.

\section{Evaluating a decider}

Accuracy on a finite test set is an estimate. The Wilson interval
\citep{wilson1927} is the standard binomial interval used here; McNemar's test
\citep{mcnemar1947} compares two classifiers on the same rows by counting
discordant pairs; \citet{dietterich1998} reviews approximate tests for classifier
comparison and recommends McNemar or resampled variants in the paired setting;
\citet{dror2018} give NLP-specific guidance and warn against tests that assume
independence across examples. \citet{card2020power} show that many NLP comparisons
are underpowered, which is why this thesis reports intervals and a replication
floor rather than stars.

Benchmark validity is a separate problem. \citet{recht2019imagenet} find that
classifier accuracy drops when the test distribution is recreated rather than
reused, \citet{ribeiro2020checklist} propose behavioural tests beyond aggregate
accuracy, and \citet{gardner2020contrast} construct contrast sets that perturb
local decision boundaries. \citet{bommasani2021foundation} situate these concerns
in the broader risk profile of foundation models. A single number on a public
test set is therefore a measurement of that set, not of the task.

This dissertation's protocol is a response to that literature: pre-register the
recipe and the decision rule, score every test row rather than a sample, report
Wilson intervals, paired bootstrap bounds and exact McNemar tests, measure a
replication floor, and keep every per-example record so that a reader can
recompute the number. The protocol is not novel as statistics; it is novel in
being applied to a decider at several candidate widths on one checkpoint.

\section{What the literature does not answer}

The gap is specific. Ranking and reranking vary candidate sets, but the candidates
come from a generator whose recall changes with the set, so accuracy and recall
are confounded. Constrained decoding varies constraints, not the trained width,
and guarantees form rather than correctness. Set-valued prediction fixes the
model and varies the set at inference, which is exactly the vacuous version of
the sweep: with the gold always present, restricting a fixed model to a sub-bank
cannot lower accuracy. Selective classification varies a threshold on one model,
not the label space the model was trained on. The public deciders publish no
width curve at all.

The experiment this thesis runs is therefore not a replication of any of them. It
trains a fresh head at each width from the same checkpoint, on the same rows, at
the same step budget, and tests at the width it trained at. The result
(Chapter~\ref{ch:results}) is that at 113.47M a 32-way interface is at or above
the full bank on three wide-bank tasks, while a narrower interface is ahead of it
by \ECostTwoToThirtyTwoMin--\ECostTwoToThirtyTwoMax{} points. That is evidence for the registered non-inferiority claim
and against the stronger claim that accuracy is flat in width. It is one seed and
one rung; the literature offers no prior estimate to compare it with, and this
dissertation does not pretend that one exists.

The same gap bounds the product claim. A clear point win on the 77-way
Banking77 cell is measured, as are clear wins on three other cited comparator
cells; a paired win on any cell is not, because no Jev number in this record
comes from the same rows. Beating Jev on ``any benchmark'' would require running
both systems on the same rows across the suite, and the stored record contains
three clear deficits --- AG News, XNLI and Language ID --- that such a run would
have to overturn rather than reinterpret. That run is listed in the remaining-work
target.
